\section{Introduction}

Machine-learning asset-pricing models can use many firm characteristics while producing a small number of latent factors or concentrated hidden representations. A selected factor count is useful for fitting and forecasting, but its interpretation as an economic dimension requires a separate argument. The distinction matters when a fitted width is described as a count of priced risks or used to fix model complexity across assets and periods.

Flexible return predictors and characteristic-based factor models summarize the cross-section in different ways \citep{gu2020empirical,kelly2019characteristics,gu2021autoencoder,chen2024deep}. Activation rank measures concentration, a factor set defines a payoff span, and a validation minimum selects a candidate under a criterion. None of these quantities alone identifies the smallest number of structural shocks. This distinction does not imply that the cited studies all make an unsupported structural claim, or that compact models lack economic value.

Related work already studies economic restrictions and interpretation. \citet{chen2026vaeer} combine variational autoencoders with economic restrictions, while \citet{zhao2026categorical} construct financially grouped factors to connect interpretation with prediction. \citet{zhang2021tradable} compare factor models through alternative HJ-distance measures. Our question concerns what an observed dimension choice establishes about its economic object. We do not directly test or refute those authors' models, and a change in ranking across different loss functions is not itself a new identification result.

We compare several pieces of evidence in one documented research sequence. The original pricing family contains six nominal dimensions and five initializations, evaluated across asset sets, weighting schemes, archived characteristic pipelines, and periods. It was retained for diagnosis after failing earlier benchmark-advancement criteria. A known-loading simulation separates evaluation-mean noise from neural optimization. Later controls isolate feature removal on a common panel, compare trained with random representations, and add a conditional autoencoder with a matched linear-loading arm. All later controls use previously studied historical data and remain exploratory.

The findings qualify each other. The archived diagnostic networks change preferred dimensions across evaluation designs. Their pipeline comparison also changes coverage, construction, and one return month, so it cannot isolate the effect of deleting six characteristics. A same-panel ablation gives positive rank correlations of 0.65--0.72 rather than reproducing the archived reversal. The locked 2020--2025 evaluation meets only one of four directional criteria, with imprecise market-attenuation estimates; legacy development had already included January 2020, preventing a research-wide untouched-holdout interpretation.

The conditional-autoencoder control produces more concentrated selection. The validation-selected $K=8$ model has development predictive $R^2$ of 0.269\% and an HJ-type distance of 5.025, compared with 5.958 for FF5 plus momentum under the same metric. These are point-estimate comparisons, not evidence of significant superiority. Within the fixed grid, $K=8$ wins 87.2\% of development resampling draws. This observation rules out a uniform description of our tested families as selection-unstable, but does not identify eight structural risks or a global optimum beyond the grid boundary.

The new control also exposes an inferential limit. Although all 15 exploratory dimension-comparison intervals include zero, a separately frozen calibration gives simultaneous coverage of 83.5\% and 87.0\%, below the nominal 95\%, in its two primary equal-distance scenarios. We therefore do not use those intervals to support equivalence, a confidence set for dimension, or structural nonidentification. This failure concerns the specified statistical procedure and simulated moments; it is not evidence for an economic mechanism.

The geometry and known-loading exercises address distinct questions. Absolute activation rank falls as layers narrow, while width-normalized rank increases. Training adds concentration relative to matched random networks, but final-layer rank does not meet the specified incremental pricing-loss prediction criterion. In the known-loading simulation, noisy evaluation reduces average exact-dimension recovery from 64.2\% to 32.6\%. A finite nested Gaussian calculation illustrates the familiar difference between prediction and consistent exact-model selection \citep{shao1993cv,yang2007cv}; it is not a general theorem for fitted neural networks.

The study's contribution is the comparison of these objects and the attribution supplied by the bounded controls, rather than a new selection theorem or a uniformly superior pricing model. The results constrain explanations based on feature deletion, architectural compression, and a single validation minimum. They also show where the present evidence stops: favorable performance and concentrated candidate selection do not establish an invariant economic dimension, while unstable selection and non-rejection do not establish its absence. The separate updating exercise is retained in the supplement as a retrospective boundary case, not as a direct test of the pricing networks.

Section~2 sets out the objects and the known-loading calculation. Section~3 describes the data, chronology, and estimators. Section~4 reports the archived comparisons, matched controls, conditional-autoencoder extension, and selection simulation. Sections~5 and 6 address the scope of updating evidence and representation geometry. Section~7 discusses the findings and limitations. \ref{app:inference-calibration} reports the new inference calibration.
